\begin{longtable}{llllll}
\caption{Kiểm định giả định của mô hình VaR} \label{tab:T2_assumption_tests} \\
\toprule
label & test & statistic & p\_value & sig & conclusion \\
\midrule
\endfirsthead
\caption[]{Kiểm định giả định của mô hình VaR} \\
\toprule
label & test & statistic & p\_value & sig & conclusion \\
\midrule
\endhead
\midrule
\multicolumn{6}{r}{Continued on next page} \\
\midrule
\endfoot
\bottomrule
\endlastfoot
VN-Index & Jarque–Bera & 1607,580 & <0,001 & *** & Bác bỏ phân phối chuẩn → Normal/MC\_Normal dễ đánh giá thấp rủi ro đuôi \\
VN-Index & Anderson–Darling & 59,083 & <0,001 & *** & Bác bỏ chuẩn \\
VN-Index & ADF & -16,715 & <0,001 & *** & Bác bỏ nghiệm đơn vị → dừng \\
VN-Index & KPSS & 0,056 & 0,100 & NaN & Không bác bỏ tính dừng \\
VN-Index & ARCH-LM(5) & 639,314 & <0,001 & *** & Có hiệu ứng ARCH (biến động theo cụm) → cần GARCH \\
VN-Index & ARCH-LM(10) & 766,062 & <0,001 & *** & Có hiệu ứng ARCH (biến động theo cụm) → cần GARCH \\
VN-Index & Ljung–Box Q(10) r & 192,593 & <0,001 & *** & Có tự tương quan trong lợi suất \\
VN-Index & Ljung–Box Q(10) r² & 2084,570 & <0,001 & *** & Có tự tương quan trong r² (cụm biến động) \\
S\&P 500 (USD) & Jarque–Bera & 29219,200 & <0,001 & *** & Bác bỏ phân phối chuẩn → Normal/MC\_Normal dễ đánh giá thấp rủi ro đuôi \\
S\&P 500 (USD) & Anderson–Darling & inf & <0,001 & *** & Bác bỏ chuẩn \\
S\&P 500 (USD) & ADF & -14,703 & <0,001 & *** & Bác bỏ nghiệm đơn vị → dừng \\
S\&P 500 (USD) & KPSS & 0,171 & 0,100 & NaN & Không bác bỏ tính dừng \\
S\&P 500 (USD) & ARCH-LM(5) & 1337,680 & <0,001 & *** & Có hiệu ứng ARCH (biến động theo cụm) → cần GARCH \\
S\&P 500 (USD) & ARCH-LM(10) & 1417,770 & <0,001 & *** & Có hiệu ứng ARCH (biến động theo cụm) → cần GARCH \\
S\&P 500 (USD) & Ljung–Box Q(10) r & 118,935 & <0,001 & *** & Có tự tương quan trong lợi suất \\
S\&P 500 (USD) & Ljung–Box Q(10) r² & 4802,300 & <0,001 & *** & Có tự tương quan trong r² (cụm biến động) \\
S\&P 500 (VND) & Jarque–Bera & 27537,200 & <0,001 & *** & Bác bỏ phân phối chuẩn → Normal/MC\_Normal dễ đánh giá thấp rủi ro đuôi \\
S\&P 500 (VND) & Anderson–Darling & inf & <0,001 & *** & Bác bỏ chuẩn \\
S\&P 500 (VND) & ADF & -14,912 & <0,001 & *** & Bác bỏ nghiệm đơn vị → dừng \\
S\&P 500 (VND) & KPSS & 0,120 & 0,100 & NaN & Không bác bỏ tính dừng \\
S\&P 500 (VND) & ARCH-LM(5) & 1298,230 & <0,001 & *** & Có hiệu ứng ARCH (biến động theo cụm) → cần GARCH \\
S\&P 500 (VND) & ARCH-LM(10) & 1383,310 & <0,001 & *** & Có hiệu ứng ARCH (biến động theo cụm) → cần GARCH \\
S\&P 500 (VND) & Ljung–Box Q(10) r & 129,303 & <0,001 & *** & Có tự tương quan trong lợi suất \\
S\&P 500 (VND) & Ljung–Box Q(10) r² & 4665,200 & <0,001 & *** & Có tự tương quan trong r² (cụm biến động) \\
USD/VND & Jarque–Bera & 9388000,000 & <0,001 & *** & Bác bỏ phân phối chuẩn → Normal/MC\_Normal dễ đánh giá thấp rủi ro đuôi \\
USD/VND & Anderson–Darling & inf & <0,001 & *** & Bác bỏ chuẩn \\
USD/VND & ADF & -25,906 & <0,001 & *** & Bác bỏ nghiệm đơn vị → dừng \\
USD/VND & KPSS & 0,225 & 0,100 & NaN & Không bác bỏ tính dừng \\
USD/VND & ARCH-LM(5) & 4,128 & 0,531 & NaN & Không thấy hiệu ứng ARCH \\
USD/VND & ARCH-LM(10) & 4,225 & 0,937 & NaN & Không thấy hiệu ứng ARCH \\
USD/VND & Ljung–Box Q(10) r & 32,868 & <0,001 & *** & Có tự tương quan trong lợi suất \\
USD/VND & Ljung–Box Q(10) r² & 4,354 & 0,930 & NaN & Không có tự tương quan trong r² (cụm biến động) \\
DM 50/50 & Jarque–Bera & 5828,940 & <0,001 & *** & Bác bỏ phân phối chuẩn → Normal/MC\_Normal dễ đánh giá thấp rủi ro đuôi \\
DM 50/50 & Anderson–Darling & 55,097 & <0,001 & *** & Bác bỏ chuẩn \\
DM 50/50 & ADF & -16,656 & <0,001 & *** & Bác bỏ nghiệm đơn vị → dừng \\
DM 50/50 & KPSS & 0,051 & 0,100 & NaN & Không bác bỏ tính dừng \\
DM 50/50 & ARCH-LM(5) & 865,531 & <0,001 & *** & Có hiệu ứng ARCH (biến động theo cụm) → cần GARCH \\
DM 50/50 & ARCH-LM(10) & 950,304 & <0,001 & *** & Có hiệu ứng ARCH (biến động theo cụm) → cần GARCH \\
DM 50/50 & Ljung–Box Q(10) r & 132,328 & <0,001 & *** & Có tự tương quan trong lợi suất \\
DM 50/50 & Ljung–Box Q(10) r² & 2881,730 & <0,001 & *** & Có tự tương quan trong r² (cụm biến động) \\
\end{longtable}
